% ===========================================================================
\section{Conclus\~ao}
\label{sec:conclusao}
% ===========================================================================

\subsection{Retomada do problema}

O trabalho partiu da seguinte pergunta: \'e poss\'ivel produzir, de forma integralmente
reproduz\'ivel e automatizada por um agente conectado via \emph{Model Context Protocol}, um
modelo tridimensional de um autom\'ovel esportivo cuja geometria e materiais sejam fi\'eis a
refer\^encias fotogr\'aficas, atendendo simultaneamente a restri\c{c}\~oes de topologia de
malha e a um envelope dimensional verificado?

A resposta que os resultados sustentam \'e \emph{sim, com ressalvas precisamente
localizadas}. O modelo existe, \'e reconstru\'ivel por um \'unico comando, tem topologia
integralmente quadrangular nos corpos principais, envelope dentro da toler\^ancia declarada
e os elementos de identidade da variante presentes e verific\'aveis. A fidelidade de
silhueta e de tom de cor, por\'em, \'e parcial, e as causas de cada desvio foram
identificadas.

\subsection{Respostas aos objetivos}

Os sete objetivos espec\'ificos foram atendidos quanto ao m\'etodo e \`a evid\^encia: as
dimens\~oes e propor\c{c}\~oes foram extra\'idas e registradas com origem declarada
(Tabela~\ref{tab:dimensoes}); o casco \'e gerado por \emph{loft} sem opera\c{c}\~oes
booleanas e sem n-gons (750 quads); os elementos de identidade da variante est\~ao
modelados como geometria real (Figuras~\ref{fig:roda-clay} e~\ref{fig:cmp-a}); a
biblioteca de materiais tem 13 materiais fisicamente plaus\'iveis; a cadeia de
automa\c{c}\~ao via MCP executa \emph{build} e renders por \emph{script}; o relat\'orio de
malha \'e emitido a cada \emph{build}; e a compara\c{c}\~ao com as refer\^encias est\'a
registrada por crit\'erio (Tabela~\ref{tab:crit}). Os crit\'erios parciais (CA-004 e
CA-007) n\~ao decorrem de objetivos n\~ao cumpridos, e sim de que a \emph{qualidade}
exigida em dois crit\'erios perceptuais n\~ao foi integralmente alcan\c{c}ada.

\subsection{Conclus\~oes}

Tr\^es conclus\~oes podem ser afirmadas com o suporte das evid\^encias apresentadas.

\textbf{Primeira: as restri\c{c}\~oes topol\'ogicas devem ser impostas na
constru\c{c}\~ao.} A aus\^encia total de n-gons e de arestas n\~ao variedade em uma cena de
32.632 faces foi obtida sem qualquer etapa de reparo. O m\'etodo de tampa por leque
dobrado e a estrutura de encadeamento por an\'eis s\~ao a raz\~ao disso. A
implica\c{c}\~ao \'e direta: a qualidade estrutural de um modelo procedimental pode ser uma
\emph{propriedade do gerador}, com garantia mais forte que a de qualquer
inspe\c{c}\~ao posterior.

\textbf{Segunda: a verifica\c{c}\~ao por medi\c{c}\~ao e a verifica\c{c}\~ao por
inspe\c{c}\~ao s\~ao complementares.} O relat\'orio num\'erico confirmou o envelope, a
topologia e a aus\^encia de artefatos estruturais, mas foi incapaz de detectar a
aus\^encia da linha de car\'ater. A inspe\c{c}\~ao visual comparativa identificou esse
desvio, mas n\~ao teria detectado um n-gon oculto ou um desvio de mil\'imetros no envelope.
Concluir que qualquer uma das duas bastaria seria um erro metodol\'ogico.

\textbf{Terceira: a media\c{c}\~ao por MCP agrega auditabilidade, com um limite
expl\'icito.} Cada opera\c{c}\~ao fica serializada como mensagem do protocolo, o que
transforma a sess\~ao de modelagem em um registro inspecion\'avel. O limite \'e que a
inspe\c{c}\~ao visual mediada por captura bidimensional \'e inferior \`a
inspe\c{c}\~ao interativa, e que o agente n\~ao substitui o julgamento de engenharia sobre
o que medir, o que aceitar e o que perseguir.

Do ponto de vista do crit\'erio final de um trabalho de engenharia --- a exist\^encia de um
artefato defens\'avel, verific\'avel e rastre\'avel --- o resultado \'e satisfat\'orio
dentro dos limites declarados. O modelo n\~ao \'e indistingu\'ivel de uma fotografia em
\emph{close} de superf\'icie; \'e, por\'em, integralmente reconstru\'ivel, medido, e cada
uma de suas fei\c{c}\~oes \'e rastre\'avel at\'e uma decis\~ao documentada.

\subsection{Trabalhos futuros}

As pend\^encias identificadas definem uma agenda de quatro itens, em ordem decrescente de
rela\c{c}\~ao custo-benef\'icio: (1)~\textbf{an\'eis de proximidade} na linha de
car\'ater e nos \emph{lips} de para-lama, que ataca diretamente CA-004 com uma t\'ecnica
j\'a validada no projeto (os an\'eis de reten\c{c}\~ao do nariz e da cauda), ou
alternativamente peso de dobra por aresta; (2)~\textbf{cadeia colorim\'etrica
controlada}, capturando uma refer\^encia com alvo de cor conhecido e aplicando o mesmo
mapeamento de tons ao render e \`a refer\^encia, para que CA-007 deixe de ser julgamento e
passe a ser medi\c{c}\~ao; (3)~\textbf{glifos de contorno fechado} nos emblemas, o que
elimina as 1.776 arestas de borda e os tri\^angulos de chanfro; e (4)~\textbf{valida\c{c}\~ao
cruzada por reconstru\c{c}\~ao densa}, empregando \emph{Structure-from-Motion} sobre as
mesmas fotografias para obter um envelope independente e atacar a principal amea\c{c}a de
validade interna. Dois desdobramentos de pesquisa completam a agenda: generalizar o gerador
para uma fam\'ilia de ve\'iculos, avaliando se a parametriza\c{c}\~ao por
esta\c{c}\~oes se sustenta quando a morfologia varia substancialmente; e estudar
sistematicamente o efeito da media\c{c}\~ao por agentes sobre a qualidade do artefato, com
avalia\c{c}\~ao por m\'ultiplos avaliadores e compara\c{c}\~ao controlada contra um fluxo
manual --- o controle experimental que este trabalho, por limita\c{c}\~ao de escopo,
n\~ao executou.

\subsection{Considera\c{c}\~oes finais}

Este trabalho come\c{c}ou com uma cena vazia e seis fotografias. Termina com um modelo de
130 objetos, cuja geometria \'e descrita por cerca de 1.200 linhas de c\'odigo que qualquer
pessoa pode ler, executar e contestar. Esse \'e, talvez, o resultado mais duradouro: n\~ao
o modelo em si, que envelhecer\'a com a vers\~ao do \emph{software}, mas a
demonstra\c{c}\~ao de que a modelagem tridimensional pode ser conduzida como
engenharia --- com especifica\c{c}\~ao, medi\c{c}\~ao, evid\^encia e registro do desvio.
